% !TeX root = ../main.tex

\section{Evaluation}\label{sec:evaluation}

We first present the evaluation setup, and then report the evaluation results~of~research~questions.

\subsection{Evaluation Setup}



\todel{We conduct extensive experiments to evaluate the effectiveness, efficiency and usefulness of \tool.}

\textbf{Research Questions.} We design our evaluation to answer the following four research questions.

\begin{itemize}[leftmargin=*]
    \item \textbf{RQ1 Effectiveness Evaluation:} How is the effectiveness~of \tool, compared with the state-of-the-art malicious package detection approaches?
    \item \textbf{RQ2 Efficiency Evaluation:} How is the performance overhead of \tool?% in offline training and online prediction?
    \item \tosemrevision{\textbf{RQ3 Dataset Scale Evaluation:} How does the scale of the dataset affect the effectiveness of \tool?}
    \item \textbf{RQ4 Ablation Study:} How behavior sequence contributes to the effectiveness of \tool?
    \item \textbf{RQ5 Usefulness Evaluation:} How useful is \tool in real-world detection on PyPI and~NPM?
\end{itemize}



\textbf{Dataset Collection.} We constructed the dataset as follows. For malicious packages, we collected malicious samples from three sources. We collected \new{438} malicious PyPI packages and \new{1,788} malicious NPM packages from Backstabber's Knife Collection~\cite{ohm2020backstabber}. We added \new{88} malicious PyPI packages from \textsc{MalOSS}'s dataset~\cite{duan2020towards}. \tosemrevision{We also sampled \new{361} packages out of a total of \tosemrevision{5,874} packages from pypi\_malregistry \cite{guo2023empirical}, with a confidence level of 95\% and a margin of error of 5\%.} For benign packages, we collected \new{2,398} benign PyPI packages from Vu et al.'s dataset~\cite{vu2023bad}. Following the procedure in prior work~\cite{ohm2022feasibility, vu2023bad}, we selected the \new{5,000} most depended upon packages in the NPM registry as the dataset of benign NPM packages. However, we successfully downloaded \new{4,993} of them. In total,~we curated a dataset of PyPI and NPM packages with \new{2,314} malicious packages and \new{7,391} benign packages, as listed in Table~\ref{table:dataset} where \textit{Mixed} denotes the summation across PyPI and NPM. \tosemrevision{We also provide the average file count and average thousand lines of code (KLOC) for each package.} 




\begin{table}[!t]
    \centering
        \small
    \caption{\tosemrevision{Statistics of the Dataset}}\label{table:dataset}
    \vspace{-10pt}
    \begin{tabular}{m{2cm}m{1.6cm}m{1.6cm}m{1.6cm}m{1.6cm}m{1.6cm}m{1.6cm}}
        \toprule
        \multirow{2}{*}{Package Registry} & \multicolumn{3}{c}{Malicious} & \multicolumn{3}{c}{Benign} \\
        \cmidrule(lr){2-4}
        \cmidrule(lr){5-7}
        & \# Packages & \tosemrevision{Aver. Files \#} & \tosemrevision{Aver. KLOC \#} & \# Packages & \tosemrevision{Aver. Files \#} & \tosemrevision{Aver. KLOC \#} \\
        \midrule
        % PyPI & 887 & 4,977 & 875,678 & 2,398 & 180,568 & 46,741,729 \\
        % NPM & 1,788 & 7,911 & 2,935,127 & 4,993 & 214,312 & 46,215,862 \\
        % Mixed & 2,675 & 12,888 & 3,810,805 & 7,391 & 394,880 & 92,957,591 \\
        PyPI & 887 & 5.6 & 0.987 & 2,398 & 75.3 & 19.492 \\
        NPM & 1,788 & 4.4 & 1.642 & 4,993 & 214.3 & 9.256 \\
        Mixed & 2,675 & 4.8 & 1.425 & 7,391 & 53.4 & 12.577 \\
        \bottomrule
    \end{tabular}
\end{table}

%We selected \todo{5,000} benign NPM packages from NPM rank\footnote{https://gist.github.com/anvaka/8e8fa57c7ee1350e3491} ranked by numbers of most depended-upon packages. We download latest five versions of those benign packages if there are any. 

\textbf{RQ Setup.} For \textbf{RQ1}, we aim to compare the effectiveness of \tool with state-of-the-art approaches. \tosemrevision{Specifically, we selected \textsc{Amalfi}~\cite{sejfia2022practical}, \textsc{SAP}~\cite{ladisa2023feasibility} and \textsc{MPHunter}~\cite{liang2023needle} as the state-of-the-art learning-based approaches}, and OSS Detect Backdoor~\cite{Ossdetectbackdoor} and Bandit4Mal~\cite{bandit4mal} as the state-of-the-art rule-based approaches. We adopted the same configuration used in the original paper~\cite{sejfia2022practical, ladisa2023feasibility, liang2023needle}. For OSS Detect Backdoor and Bandit4Mal, we set the threshold of 3 alerts to distinguish malicious and benign packages, which was observed as the optimal threshold by Vu et al.~\cite{vu2023bad}.~We did not compare with \textsc{MalOSS}~\cite{duan2020towards} as it used dynamic features and thus failed to run on many packages. We also did~not compare with Malware Checks~\cite{malwarechecks} as the PyPI team removed it recently. We split the dataset into training and testing dataset by 9:1. We measured the precision and recall of all the approaches with 10-fold cross-validation on the testing datasets. We compared these approaches in three detection scenarios, i.e., mono-lingual scenario (i.e., train and test on the same ecosystem), cross-lingual scenario (i.e., train on one ecosystem and test on the other), and bi-lingual scenario (i.e., train on two ecosystems and test on one of the two ecosystems). \tosemrevision{We trained \tool with BERT~\cite{kenton2019bert}, RoBERTa~\cite{liu2019roberta} and T5~\cite{raffel2020exploring} on the PyPI, NPM and Mixed training datasets, denoted by \tool$_{BERT}$, \tool$_{RoBERTa}$ and \tool$_{T5}$. We also trained \textsc{Amalfi} with decision tree (\textsc{Amalfi}$_{DT}$), naive bayes (\textsc{Amalfi}$_{NB}$) and SVM (\textsc{Amalfi}$_{SVM}$) on the same training datasets. We only evaluate \textsc{SAP} in both mono-lingual and bi-lingual scenarios, and \textsc{MPHunter} in the mono-lingual scenario for PyPI packages, as their do not support adaptation to other scenarios.}



For \textbf{RQ2}, we measured the time overhead of \tool~in offline training and online prediction. \tosemrevision{For \textbf{RQ3}, we evaluated the effectiveness of \tool across different dataset scales. We incrementally increased the proportion of our training dataset from 10\% to 100\% in 10\% steps. We trained \tool and \tosemrevision{\textsc{Amalfi}$_{DT}$} in the bi-lingual scenario using these varying dataset proportions. The trained models were then tested using the same testing dataset across all scales. The \tosemrevision{average F1-Score over 10-fold} cross-validation was used to measure overall effectiveness.}


For \textbf{RQ4}, we created three ablated versions of \tool and compared their effectiveness with the original version of \tool using the same dataset in \textbf{RQ1}. Specifically, to evaluate~the~impact of sequential behavior, we created “\tool w/o Seq” by removing behavior sequence generator and feeding textual descriptions directly into maliciousness classifier. To evaluate the impact of textual description transformation, we created ``\tool w/o Text'' by removing textual description transformation in maliciousness classifier. To evaluate the impact of \tosemrevision{pre-trained language models}, we created ``\tool w/ DT'' by removing behavior sequence generator and feeding features as a vector into~the~decision~tree.

For \textbf{RQ5}, we ran \tosemrevision{\tool using the BERT model} against the newly published packages in PyPI and NPM for \todel{over \todo{6} weeks}\new{over nine months}.  \new{This monitoring process started \tosemminorrevision{on} March 03 2023 for PyPI and April 01 2023 for NPM, and ended at October~30~2023}\todel{May~27~2023}. \tool analyzed \new{599,493} PyPI package versions and \new{324,145} NPM package versions. We manually confirmed the potentially malicious package versions flagged~by~\tool.\tosemrm{, and reported them to the official PyPI and NPM teams.} \tosem{Specifically, two authors manually assessed the maliciousness of the packages in the dimension of static characteristics, dynamic behaviors and third-party detection tools. Two authors conducted a manual assessment to determine the maliciousness of the packages. This process covered comprehending the source code, inspecting the dynamic behaviors, and referencing detection results from third-party detection tools (e.g., virustotal~\cite{virustotal}). The two authors are postgraduate and undergraduate students with security/malware background. The human inspection was conducted daily, costing one hour each day, over a period of ten months. The average number of reviewed packages daily was 30. All packages confirmed by two authors as malicious were reported to the official PyPI and NPM teams, all of which were accepted by the two registries. }Further,~by adding the confirmed~malicious and benign packages into the original training datasets (i.e., incremental learning),~we~retrained \tosemrevision{\tool using the BERT model} to evaluate whether \tool could learn from new data to improve its effectiveness. We also analyzed the malicious intentions and triggering scenarios of the new malicious packages.

\begin{equation}\label{eq:pre_rec}
    \begin{small}
    \begin{aligned}
        \vspace{-5pt}
        Pre. = \frac{\mid  MP_{tool} \cap MP_{test} \mid}{\mid MP_{tool} \mid }\\
        %    \mathrm{,~} 
        Rec. =  \frac{\mid  MP_{tool} \cap MP_{test}  \mid }{\mid MP_{test} \mid}\\
        % 文本表示F1-Score
        \tosemrevision{F1-Score = \frac{2 \times Pre. \times Rec.}{Pre. + Rec.}}
    \end{aligned}
    \end{small}
\end{equation}

\tosem{\textbf{Evaluation Metric.} The precision (Pre.), recall (Rec.) used in \textbf{RQ1}, \textbf{RQ4}, \textbf{RQ5} and \tosemrevision{F1-Score used in \textbf{RQ3}} are defined in Equation~\ref{eq:pre_rec}. Specifically, $MP_{tool}$ denotes the set of malicious packages predicted by tools (e.g., \tool), and $MP_{test}$ denotes the set of true malicious packages in the testing dataset. Precision is calculated by the proportion of true positives (i.e., $MP_{tool} \cap MP_{test}$) in the predicted malicious packages (i.e., $MP_{tool}$), and recall is calculated by the proportion of true positives (i.e., $MP_{tool} \cap MP_{test}$) in the ground truth (i.e., $MP_{test}$).}



\begin{table}[!t]
    \centering
    \footnotesize
    \caption{\tosemrevision{Evaluation Result in Mono-Lingual Scenario for Learning-based Approaches}}\label{table:result-mono-lingual}
    \vspace{-10pt}
    \begin{tabular}{m{0.40cm}m{0.40cm}m{0.46cm}m{0.46cm}m{0.46cm}m{0.46cm}m{0.46cm}m{0.46cm}m{0.46cm}m{0.46cm}m{0.46cm}m{0.46cm}m{0.51cm}m{0.51cm}m{0.46cm}m{0.46cm}m{0.46cm}m{0.46cm}}   
        \noalign{\hrule height 1pt}
        \multirow{2}{*}{Train}  & \multirow{2}{*}{Test} & \multicolumn{2}{c}{\shortstack[l]{\tool$_{BERT}$}} & \multicolumn{2}{c}{\tosemrevision{\shortstack[l]{\tool$_{RoBERTa}$}}} & \multicolumn{2}{c}{\tosemrevision{\shortstack[l]{\tool$_{T5}$}}}  & \multicolumn{2}{c}{\shortstack[l]{\textsc{Amalfi}$_{DT}$}} & \multicolumn{2}{c}{\shortstack[l]{\textsc{Amalfi}$_{NB}$}} & \multicolumn{2}{c}{\shortstack[l]{\textsc{Amalfi}$_{SVM}$}} & \multicolumn{2}{c}{\tosemrevision{\textsc{SAP}}} & \multicolumn{2}{c}{\tosemrevision{\textsc{MPHunter}}}\\
    \cmidrule(lr){3-4}
    \cmidrule(lr){5-6}
    \cmidrule(lr){7-8}
    \cmidrule(lr){9-10}
    \cmidrule(lr){11-12}
    \cmidrule(lr){13-14}
    \cmidrule(lr){15-16}
    \cmidrule(lr){17-18}
    & & Pre. & Rec. & Pre. & Rec. & Pre. & Rec. & Pre. & Rec. & Pre. & Rec. & Pre. & Rec. & Pre. & Rec. & Pre. & Rec.\\ 
    \noalign{\hrule height 1pt}
    PyPI & PyPI & 95.8\% & 89.4\% & \tosemrevision{96.0\%} & \tosemrevision{91.7\%} & \tosemrevision{93.4\%} & \tosemrevision{70.9\%} & 81.9\% & 83.9\% & 59.5\% & 7.5\% & 30.3\% & 30.9\% & \tosemrevision{83.9\%} & \tosemrevision{61.5\%} & \tosemrevision{21.8\%} & \tosemrevision{85.4\%}\\
    NPM & NPM & 98.2\% & 91.8\% & \tosemrevision{98.5\%} & \tosemrevision{92.9\%} & \tosemrevision{98.9\%} & \tosemrevision{91.0\%} & 92.7\% & 86.0\% & 92.1\% & 77.1\% & 10.2\% & 26.7\%  & \tosemrevision{90.5\%} & \tosemrevision{35.4\%} & \tosemrevision{--} & \tosemrevision{--}\\
    \noalign{\hrule height 1pt}
    \end{tabular}
\end{table}

\begin{table}[!t]
    \centering
    \footnotesize
    \caption{Evaluation Result in Mono-Lingual Scenario for Rule-Based Approaches}\label{table:result-rule-based}
    \vspace{-10pt}
    \begin{tabular}{ccccc}
        \noalign{\hrule height 1pt}
        \multirow{2}{*}{Test} & \multicolumn{2}{c}{OSS Detect Backdoor} & \multicolumn{2}{c}{Bandit4Mal} \\
        \cmidrule(lr){2-3}
        \cmidrule(lr){4-5}
         & Pre. & Rec. & Pre. & Rec. \\ 
        \noalign{\hrule height 1pt}
        PyPI & 19.9\% & 60.4\% & 24.9\% & 74.4\%\\
        NPM & 41.1\% & 85.1\%  &  -- & -- \\
        \noalign{\hrule height 1pt}
    \end{tabular}
\end{table}


\subsection{Effectiveness Evaluation (RQ1)}

\todel{We report the effectiveness evaluation results with respect~to the mono-lingual, cross-lingual and bi-lingual scenarios.}

\textbf{Mono-Lingual Scenario.} The result of learning-based approaches in mono-lingual scenario is reported in Table~\ref{table:result-mono-lingual} and the result of rule-based approaches is reported in Table~\ref{table:result-rule-based}. \tosemrevision{Overall, \tool$_{RoBERTa}$ outperforms all the other approaches. In PyPI, \tool$_{RoBERTa}$ achieves a precision of 96.0\% and a recall of 91.7\%. It achieves the best result, with slight advantages compared with \tool$_{BERT}$ and \tool$_{T5}$. It outperforms the state-of-the-art \textsc{SAP} by 12.1\% in precision, and \textsc{MPHunter} by 6.3\% in recall. In NPM, \tool$_{RoBERTa}$ achieves a precision of 98.5\% and a recall of 92.9\%. It outperforms \tool$_{T5}$ with a 1.9\% higher recall and a 0.4\% lower precision, and \tool$_{BERT}$ with a 0.3\% higher precision and 1.1\% higher recall. It outperforms the state-of-the-art \textsc{Amalfi}$_{DT}$ by 5.8\% in precision and 6.9\% in recall.} The two rule-based approaches have low precision and recall. 

\textbf{Cross-Lingual Scenario.} The result of cross-lingual scenario is presented in Table~\ref{table:result-cross-lingual}. 
\tosemrevision{When the model is trained on NPM packages and tested on PyPI packages, \tool$_{RoBERTa}$~obtains a highest precision of 75.1\%. However, \tool$_{T5}$ achieves a more balanced performance with a precision of 65.2\% and a recall of 63.4\%. Although \textsc{Amalfi}${SVM}$ has better recall, \tool$_{T5}$ significantly surpasses it in precision, resulting in an overall F1-Score advantage of 13.3\%. When the model~is~trained~on PyPI packages and tested on NPM packages, \tool$_{RoBERTa}$ obtains a highest precision of 46.5\%, and a recall of 92.7\%. Although \textsc{Amalfi}${SVM}$ has better recall, \tool$_{RoBERTa}$ significantly surpasses it in precision, resulting in a F1-Score advantage of 17.7\%.}
This result potentially owes to our well-abstracted feature set. Notice that \tool in the cross-lingual scenario achieves a lower precision and recall than in the mono-lingual scenario, which indicates that some malicious behaviors might be not common, and thus bi-lingual knowledge fusing is necessary.



\textbf{Bi-Lingual Scenario.} The result of bi-lingual scenario is shown in Table~\ref{table:result-bi-lingual}. 
\tosemrevision{Overall, \tool$_{RoBERTa}$ outperforms all the other approaches in PyPI and NPM packages. For testing PyPI packages using the model trained on mixed packages, \tool$_{BERT}$ obtains a precision of 95.0\% and a recall of 92.0\%, and \tool$_{RoBERTa}$ obtains a precision of 96.1\% and a recall of 90.9\%. They achieve very similar results, with \tool$_{RoBERTa}$ outperforming \tool$_{BERT}$ by 1.1\% in precision, while \tool$_{BERT}$ surpasses \tool$_{RoBERTa}$ by the same margin in recall. Comparing with the state-of-the-art, \tool$_{RoBERTa}$ outperform \textsc{Amalfi}$_{DT}$ by 12.9\% in precision and 9.8\% in recall. For testing NPM packages using the model trained on mixed packages, \tool$_{RoBERTa}$ achieves a precision of 98.9\% and a recall of 93.9\%. It outperforms \tool$_{BERT}$ by 0.2\% in precision and 0.9\% in recall, and surpasses \tool$_{T5}$ by 1.9\% in recall. Comparing with the state-of-the-art, \tool$_{RoBERTa}$ outperforms \textsc{Amalfi}$_{DT}$ by 6.8\% in precision 8.1\% in recall. Although \textsc{Amalfi}$_{SVM}$ outperforms \tool in recall by 0.7\%, it suffers from a very low precision of 28.8\%.}



\textbf{Comparing Bi-lingual with Mono-lingual.} \tosemrevision{Comparing the effectiveness of \tool in bi-lingual and mono-lingual scenarios, we observe that \tool$_{BERT}$, \tool$_{RoBERTa}$, and \tool$_{T5}$ trained on mixed packages outperform their mono-lingual counterparts. For detecting malicious PyPI packages, the bi-lingual model of \tool achieves an average precision increase of 0.2\% and an average recall increase of 0.2\%. For NPM packages, the bi-lingual model achieves an average precision increase of 0.3\% and an average recall increase of 1.1\%. These results indicate that our bi-lingual knowledge fusion is useful.}

\tosemminorrevision{When comparing the performance of \tool$_{BERT}$ and \tool$_{RoBERTa}$, we observe that \tool$_{RoBERTa}$ generally outperforms \tool$_{BERT}$ in most scenarios. This advantage can be attributed to RoBERTa’s use of more advanced training strategies, including dynamic masking, larger batch sizes, and an increased number of training steps, which enable the model to learn more efficiently. However, in cross-lingual and bi-lingual tasks, \tool$_{BERT}$ achieves higher or comparable recall compared to \tool$_{RoBERTa}$. Further analysis reveals that for malicious packages detected by \tool$_{BERT}$ but missed by \tool$_{RoBERTa}$, the sequences tend to be longer. This suggests that \tool$_{RoBERTa}$ may be less effective at capturing long-range dependencies within longer sequences, which are essential for recognizing complex malicious patterns. A potential explanation for this discrepancy lies in the pre-training tasks of the two models: RoBERTa does not include the Next Sentence Prediction (NSP) task during its training, while BERT does. The inclusion of NSP in BERT's training may enhance its ability to comprehend long-context relationships, making it better at detecting certain malicious packages that RoBERTa misses.}

% 14.1 7.8
% 5.8 6.9

\textbf{Summary.} \tosemrevision{\tool$_{RoBERTa}$ outperforms the state-of-the-art by averagely 10.0\% in precision and 7.4\% in recall in the mono-lingual scenario, and averagely 9.9\% in precision and 8.9\% in recall in the bi-lingual scenario. \tool learns bi-lingual knowledge from two package registries, which contributes to an average increase of 0.3\% in precision and 0.7\% in recall.}

% \tool learns bi-lingual knowledge from two package registries, which contributes to an average increase of 2.8\% in precision and 6.3\% in recall. 
% \tosemrevisiontodo{In the mono-lingual scenario, \tool outperforms the state-of-the-art by in the mono-lingual scenario, by in the cross-lingual scenario, and by in the bi-lingual scenario. }

% \textbf{\textit{Summary.}} \new{\tool outperforms the state-of-the-art by \new{8.6\%}\todel{\todo{10.2\%}} in precision with a trivial loss of 0.2\% in recall\todel{\todel{\todo{3.7\%}} in recall} in the mono-lingual~scenario, by \todo{35.6\%} in precision and \todo{73.4\%} in recall in the cross-lingual scenario, and by \todo{10.8\%} in precision and \todo{7.3\%} in recall in the bi-lingual scenario. \tool learns bi-lingual knowledge from two package registries, which contributes to an average increase of \todo{2.8\%} in precision and \todo{6.3\%} in recall.}




\begin{table}[!t]
    \centering
    \footnotesize
    \caption{\tosemrevision{Evaluation Result in Cross-Lingual Scenario}}\label{table:result-cross-lingual}
    \vspace{-10pt}
    \begin{tabular}{m{0.4cm}m{0.34cm}m{0.46cm}m{0.46cm}m{0.46cm}m{0.46cm}m{0.46cm}m{0.46cm}m{0.46cm}m{0.46cm}m{0.46cm}m{0.46cm}m{0.51cm}m{0.51cm}m{0.46cm}m{0.46cm}m{0.46cm}m{0.46cm}}   
        \noalign{\hrule height 1pt}
        % \multirow{2}{*}{Train}  & \multirow{2}{*}{Test}  & \multicolumn{2}{c}{\tool} & \multicolumn{2}{c}{\textsc{Amalfi-DT}} & \multicolumn{2}{c}{Amalfi-NB} & \multicolumn{2}{c}{Amalfi-SVM}\\
        \multirow{2}{*}{Train}  & \multirow{2}{*}{Test} & \multicolumn{2}{c}{\shortstack[l]{\tool$_{BERT}$}} & \multicolumn{2}{c}{\tosemrevision{\shortstack[l]{\tool$_{RoBERTa}$}}} & \multicolumn{2}{c}{\tosemrevision{\shortstack[l]{\tool$_{T5}$}}}   & \multicolumn{2}{c}{\shortstack[l]{\textsc{Amalfi}$_{DT}$}} & \multicolumn{2}{c}{\shortstack[l]{\textsc{Amalfi}$_{NB}$}} & \multicolumn{2}{c}{\shortstack[l]{\textsc{Amalfi}$_{SVM}$}}\\
    \cmidrule(lr){3-4}
    \cmidrule(lr){5-6}
    \cmidrule(lr){7-8}
    \cmidrule(lr){9-10}
    \cmidrule(lr){11-12}
    \cmidrule(lr){13-14}
    & & Pre. & Rec. & Pre. & Rec. & Pre. & Rec. & Pre. & Rec. & Pre. & Rec. & Pre. & Rec. \\ 
    \noalign{\hrule height 1pt}
    NPM & PyPI & 72.0\% & 49.1\% & \tosemrevision{75.1\%} & \tosemrevision{44.1\%} & \tosemrevision{65.2\%} & \tosemrevision{63.4\%} & 72.0\% & 18.0\% & 42.0\% & 31.5\% & 38.4\% & 76.0\% \\
    PyPI & NPM & 41.9\% & 93.0\% & \tosemrevision{46.5\%} & \tosemrevision{92.7\%} & \tosemrevision{37.5\%} & \tosemrevision{90.9\%} & 14.4\% & 4.7\% & 48.8\% & 5.2\% & 28.8\% & 94.6\% \\
    \noalign{\hrule height 1pt}
    \end{tabular}
\end{table}

\begin{table}[!t]
    \centering
    \footnotesize
    \caption{\tosemrevision{Evaluation Result in Bi-Lingual Scenario}}\label{table:result-bi-lingual}
    \vspace{-10pt}
    \begin{tabular}{m{0.45cm}m{0.40cm}m{0.46cm}m{0.46cm}m{0.46cm}m{0.46cm}m{0.46cm}m{0.46cm}m{0.46cm}m{0.46cm}m{0.46cm}m{0.46cm}m{0.51cm}m{0.51cm}m{0.46cm}m{0.46cm}}   
        \noalign{\hrule height 1pt}
        % \multirow{2}{*}{Train}  & \multirow{2}{*}{Test}  & \multicolumn{2}{c}{\tool} & \multicolumn{2}{c}{\textsc{Amalfi-DT}} & \multicolumn{2}{c}{Amalfi-NB} & \multicolumn{2}{c}{Amalfi-SVM}\\
        \multirow{2}{*}{Train}  & \multirow{2}{*}{Test}  & \multicolumn{2}{c}{\shortstack[l]{\tool$_{BERT}$}} & \multicolumn{2}{c}{\tosemrevision{\shortstack[l]{\tool$_{RoBERTa}$}}} & \multicolumn{2}{c}{\tosemrevision{\shortstack[l]{\tool$_{T5}$}}} & \multicolumn{2}{c}{\shortstack[l]{\textsc{Amalfi}$_{DT}$}} & \multicolumn{2}{c}{\shortstack[l]{\textsc{Amalfi}$_{NB}$}} & \multicolumn{2}{c}{\shortstack[l]{\textsc{Amalfi}$_{SVM}$}} & \multicolumn{2}{c}{\tosemrevision{\textsc{SAP}}}\\
    \cmidrule(lr){3-4}
    \cmidrule(lr){5-6}
    \cmidrule(lr){7-8}
    \cmidrule(lr){9-10}
    \cmidrule(lr){11-12}
    \cmidrule(lr){13-14}
    \cmidrule(lr){15-16}
    & & Pre. & Rec. & Pre. & Rec. & Pre. & Rec. & Pre. & Rec. & Pre. & Rec. & Pre. & Rec. & Pre. & Rec. \\ 
    \noalign{\hrule height 1pt}
    \multirow{2}{*}{Mixed}  & PyPI & 95.0\% & 92.0\% & \tosemrevision{96.1\%} & \tosemrevision{90.9\%} & \tosemrevision{94.6\%} & \tosemrevision{69.6\%} & 83.2\% & 81.1\% & 42.0\% & 35.6\% & 45.1\% & 37.2\% & \tosemrevision{80.2\%} & \tosemrevision{60.2\%} \\
     & NPM & 98.7\% & 93.0\% & \tosemrevision{98.9\%} & \tosemrevision{93.9\%} & \tosemrevision{98.9\%} & \tosemrevision{92.0\%} & 92.1\% & 85.8\% & 93.9\% & 76.6\% & 28.8\% & 94.6\% & \tosemrevision{90.4\%} & \tosemrevision{39.5\%} \\
    \noalign{\hrule height 1pt}
    \end{tabular}
\end{table}
